\section{Authoritative Marian doctrine beyond the four dogmas}

Mary's lifelong freedom from personal sin is authoritative doctrine closely
joined to her Immaculate Conception but not identical with its defined object.
Her spiritual motherhood and maternal cooperation in the order of grace are
taught by \emph{Lumen gentium} 56 and 60--62 and the Catechism. Her Queenship
expresses a created, limited, analogous participation in Christ's kingship.
``Mother of the Church'' joins her maternity of Christ with her relation to
his members.

These teachings have real weight without becoming additional dogmas by
editorial promotion. Their exact authority should be established from the
competent act. ``Authoritative doctrine'' is not a euphemism for optional
poetry. Neither is every sentence surrounding such teaching irreformable.

Christ alone is Redeemer and mediator in the strict sense. Mary's cooperation
is created, graced, maternal, subordinate, participated, and wholly dependent
on him. The Dicastery for the Doctrine of the Faith's \emph{Mater Populi
Fidelis} governs current constructive use: ``Co-redemptrix'' is judged
inappropriate, and ``Mediatrix'' or ``Advocate'' requires the strict
Christological limits articulated there. Historical papal or saintly use
remains evidence of use at a date, not a cumulative definition.

